\section*{Chapter \ref{ch:s1:intraday-machine-learned-alphas} --- Intraday Machine-Learned Alphas}

\begin{solution}{exo:s1:intraday-machine-learned-alphas:1}
The correlation is $\sqrt{0.25} = 0.5$, so a well-calibrated forecast has a standard deviation of $0.5 \times 0.69 = 0.34$ ticks: most forecasts are a third of a
tick or less, well below a one-tick spread.
\end{solution}

\begin{solution}{exo:s1:intraday-machine-learned-alphas:2}
The two targets share 29 of their 30 seconds, so a model tested on second $t + 1$ after training on second $t$ has seen almost all of the test target: the
test measures memory, not prediction. Training and test must be separated by at least the horizon (purging and an embargo), here by whole sessions.
\end{solution}

\begin{solution}{exo:s1:intraday-machine-learned-alphas:3}
One spread: a forecast of at least a tick in the direction traded, before fees.
\end{solution}

\begin{solution}{exo:s1:intraday-machine-learned-alphas:4}
A passive order earns the spread when it is filled; the forecast chooses the side that the price is more likely to move toward, so fills are less adversely
selected than random fills. The fill model is optimistic: it fills the order as soon as the market trades through its price, as if it were at the front of
the queue.
\end{solution}

\begin{solution}{exo:s1:intraday-machine-learned-alphas:5}
It fitted noise: with little signal at two minutes, the trees found patterns in the training sessions that do not recur, and the forecast adds variance
without predictive power, so the out-of-sample R-squared is below zero.
\end{solution}

\begin{solution}{exo:s1:intraday-machine-learned-alphas:6}
The same order flow moves the price more when the book is thin: dividing by depth makes the feature comparable across times and instruments and closer to
stationary, as Cont, Kukanov and Stoikov's linear relation with a slope inversely proportional to depth suggests.
\end{solution}

\begin{solution}{exo:s1:intraday-machine-learned-alphas:7}
Without the mid's past changes the five-second R-squared is 18.6\% for ridge (against 21.1\%) and 24.4\% for the trees (against 24.6\%): almost nothing is lost.
The order-flow and imbalance features already contain the short-term trend, which on the synthetic tape is driven by the flow itself.
\end{solution}

\begin{solution}{exo:s1:intraday-machine-learned-alphas:8}
Direction accuracy says nothing about the size of the moves against the spread: a one-second forecast is worth a fraction of a tick and a market order pays
half a spread in and half out. It must be traded passively or not at all, and the accuracy must be measured out of sample with realistic latency.
\end{solution}

\begin{solution}{pb:s1:intraday-machine-learned-alphas:1}
\begin{enumerate}
\item A seconds-to-hours forecast built from market data and traded within the day; the interval over which the move is forecast.
\item The mid's change over 5, 30 and 120 seconds, every second, in twenty one-hour sessions; a spread of about 1.07 ticks.
\item Imbalance, spread, order-flow imbalance over 1, 5 and 30 seconds, signed volume over 5 and 30, own returns over 5 and 30, microprice distance;
features must not change when the future is removed.
\item Order-flow inputs beat raw book states; the effective horizon is about two average price changes.
\item Ridge regression and sixty depth-two boosted trees.
\item Training on sessions 1--14, testing on 15--20, so overlapping targets cannot leak.
\item 21.1\% and 24.6\% at 5 seconds, 5.1\% at 30, below zero at 120; ICs 0.33--0.34, 0.23--0.28, 0.07.
\item The synthetic tape's mids trend over a few seconds, driven by its order flow.
\item Choosing the order type and price from the forecast: cross only when the forecast beats the spread.
\item Aggressive $-0.60$, $-0.48$, $-0.87$; passive 0.37, 0.33, $-0.02$; coupled 0.33, 0.22, $-0.24$ ticks a share.
\item The spread they pay exceeds the moves they predict.
\item Fills as soon as the market trades through the order's price: front of the queue.
\item \textbf{Named result.} Out-of-sample R-squared of 24.6\%, 5.1\% and $-1.5\%$ at 5, 30 and 120 seconds (boosted trees); after the spread, aggressive trading
of every forecast loses 0.60 ticks a share at five seconds while passive entries earn 0.37 and the coupled rule 0.33.
\item A single model trained on all stocks predicts better than stock-specific ones, even out of sample in stocks it has not seen.
\item Normalise features and targets by each instrument's tick, spread and depth, and train one model on all of them.
\item Passive fills would be fewer and more adversely selected; the passive and coupled P\&L would fall.
\item The forecast decays within seconds, so a delay in data or orders eats the part that is predictable.
\item Alpha-driven execution, which saves costs on trades the firm makes anyway.
\item Small: a fraction of the size at the touch, in each name.
\item What it saves or earns after the spread, in the execution it is designed with.
\end{enumerate}
\end{solution}

\begin{solution}{iq:s1:intraday-machine-learned-alphas:1}
Top-of-book and depth imbalance, order-flow imbalance over several windows, signed trade volume, the spread, the microprice's distance from the mid, recent
returns, and the same for related instruments; all normalised by depth or volatility.
\end{solution}

\begin{solution}{iq:s1:intraday-machine-learned-alphas:2}
Split in time with a gap at least as long as the horizon (purging) and a further embargo, or by whole sessions; never shuffle; report walk-forward results.
\end{solution}

\begin{solution}{iq:s1:intraday-machine-learned-alphas:3}
No: use it to decide where to rest passive orders and when not to, which earns or saves part of the spread; it becomes useless only if even passive fills
are adversely selected beyond what the forecast recovers.
\end{solution}

\begin{solution}{iq:s1:intraday-machine-learned-alphas:4}
Flatten the trees into arrays of thresholds and leaf values, update features incrementally from each message, keep everything in cache, and evaluate in
compiled code (C++ or Rust) without allocation; or precompute lookup tables for the features' quantiles.
\end{solution}

\begin{solution}{iq:s1:intraday-machine-learned-alphas:5}
Because it is traded thousands of times a day in many instruments: a small edge per trade, if it survives costs, adds up; and returns at short horizons are
mostly noise, so a 1\% R-squared can mean a correlation of 0.1, a large information coefficient.
\end{solution}

\begin{solution}{iq:s1:intraday-machine-learned-alphas:6}
Given $f$, $E[y \mid f] = \rho(\sigma_y/\sigma_f)f$. Trading $\operatorname{sign}(f)$ when $|f| > c$ earns $E[\rho(\sigma_y/\sigma_f)|f| - c \mid |f| > c]$ per
trade; with $f$ normal, $E[|f| \mid |f| > c] = \sigma_f\,\phi(c/\sigma_f)/(1 - \Phi(c/\sigma_f))$, so the P\&L grows linearly with
$\rho$ and is positive only when $\rho\sigma_y$ times that tail mean exceeds $c$: a weak correlation needs a high threshold, and trades rarely.
\end{solution}
